\section{Related Work}
\label{sec:related}
\paragraph{LLMs in game-theoretic environments} Understanding strategic behaviors of LLMs entails important societal ramifications, as online users increasingly rely on intelligent assistants to interact with other agents, potentially also LLMs. To characterize LLM's behaviors, prior studies \citep{gemp2024states,sreedhar2024simulating,de2023emergent,wu2024shall,akata2023playing,fan2024can,mao2023alympics,guo2023gpt} often adopt game theory, a mathematical framework that models cooperative behaviors of humans. These analyses involve comparing qualitative behaviors of LLM interactions against stylized entities, such as pareto optimal solutions and subgame-perfect equilibria \citep{fudenberg1991game}. LLM research in game-theoretic environments belongs to a growing body of work on multi-agent LLMs \citep{li2023theory} and their evaluation \citep{huang2024far, duan2024gtbench}. In particular, AvalonBench \citep{wang2023avalon} serves as a valuable platform for developing new multi-agent strategies. \citep{guo2023gpt,park2024ai} observe LLMs to perform in games driven by \textit{self-interest}, but falter in those that require \textit{coordination}, a behavior that emerges under altruistic and/or submissive personalities \citep{akata2023playing,guo2023gpt}. \citep{lore2023strategic,coda2024cogbench} observe different LLM families to exhibit varying levels of risk tolerance. Deferring details to experiments, we identify another source of brittleness as LLMs behave poorly when presented with numerically perturbed payoffs, even if they do not alter qualitative solutions to the game of interest. In short, there lacks a system that elicits \textit{optimal} behaviors for LLMs in game-theoretic settings.

\paragraph{Enhancing LLMs to solve games} \citep{gemp2024states} is a representative approach that aims to elicit game-solving capabilities in LLMs. It models natural dialogues as incomplete-information games, and synthesizes optimal actions by instructing an LLM to respond with specific personalities. \citep{guo2024can} proposes an LLM-based self-play algorithm that emulates Monte-Carlo Tree Search to solve zero-sum games. Broadly, these methods belong to a family of prompting strategies \cite{wei2022chain,yao2023tree,liu2024dellma} that tackle decision making instances. Our method is no exception, but differs as we imbue classical game theory in LLMs to allow for fine-grained control and analyses at each information state.

% What's lacking: decision-making / game theoretic environments with imperfect information. Prior works in (self-play) reinforcement learning that solve such tasks. Imperfect information + multi-step game
% 
% General Game -> Multiagent Game -> Imperfect/Incomplete Information Game
% Multi-agent systems / dialogues
% Reinforcement learning 

\paragraph{Game-theoretic testbeds} Stylized games such as \textit{Battle of Sexes}, \textit{Prisoner's Dilemma}, \textit{Rock-Paper-Scissors}, \textit{Stag Hunt}, and \textit{Ultimatum Game} have been extensively analyzed in the context of multi-agent systems \citep{sreedhar2024simulating,akata2023playing,fan2024can,lore2023strategic}. They represent a minimal setting characterized by small action spaces, limited number of terms, and the existence of analytical equilibria; yet they cannot capture complex interactions of real-world, multi-agent dialogues. Games that emulate real interactions are more challenging to analyze, but practically useful. Exemplar games in this category include \textit{Deal-or-No-Deal} \citep{lewis2017deal}, \textit{Multi-Round Auction} \citep{mao2023alympics}, \textit{Schedule-a-Meeting}, \textit{Trading Fruit}, \textit{Public Debate} \citep{gemp2024states}, \textit{Avalon} \citep{wang2023avalon}, \textit{Pokémon} \citep{hu2024pok}, \textit{Chess} \citep{guo2024can} and \textit{Bargaining} \citep{xia2024measuring}. And there are efforts to collect multiple games and evaluate them \citep{duan2024gtbench}. They often involve intractable action spaces (\textit{e.g.} natural language) and do not attain analytical solutions; but they are nevertheless endowed with a well-defined payoff function for practitioners to analyze the optimality of their strategy, albeit in an end-to-end manner. 

\paragraph{Workflow-aided LLM-based agent}
LLMs have been extensively utilized to decompose user requests and tasks, formulate plans, and employ various tools to execute actions \cite{ge2023openagi, wu2023autogen, li2023camel, langc, topsakal2023creating, xi2023rise}. This reliance on the innate capabilities of LLMs has propelled advancements in domains such as natural language understanding, automated reasoning, and task automation. However, depending solely on the autonomous abilities of LLMs has revealed several notable limitations, including suboptimal performance \cite{xie2024travelplanner, xia2024agentless}, lack of reliability due to output randomness \cite{inglecomprehensive, li2024survey, schwartz2023enhancing}, and the propagation of errors across sequential reasoning steps \cite{yu2023thought}.

To address these challenges, the concept of incorporating agentic workflows \cite{wu2024stateflow, zeng2023flowmind, xiao2024flowbench, li2024autoflow, li2023metaagents} into LLM-based agents has emerged. Rather than allowing LLMs to independently decompose tasks and plan actions, workflows leverage human expertise and established knowledge frameworks to guide the direction and planning processes of the LLMs. Workflow-based agents have great potential to achieve high performance and adaptability in game scenarios \citep{xi2023rise}. There have been multiple agents presented for real-world tasks \citep{jimenez2023swe, yang2024swe}, embodied agents that interact with the environment \citep{wang2023voyager, zhu2023ghost}, agents that can learn from playing a text-based game \citep{xu2023language}, and agent operating on games with imperfect information \citep{guo2023suspicion}.

\paragraph{Game-theoretic Workflow in this Paper} In this paper, we integrate game-theoretic principles into the reasoning processes of LLMs prior to decision-making. By guiding the models to derive rational strategies and make decisions based on these strategies, we aim to enhance their ability to perform effectively in strategic settings.

To the best of our knowledge, this work is the first to combine LLM-based agents with structured workflows inspired by game theory to enhance strategic decision-making capabilities. We address both complete-information and incomplete-information settings, drawing upon classical game-theoretic principles to guide the reasoning processes of LLMs. Specifically, we introduce a novel algorithm for negotiation under incomplete information, enabling agents to perform Bayesian updates and make rational decisions based on limited knowledge of other players. This approach bridges the gap between traditional game theory and LLM-based agents, advancing the applicability of LLMs in complex strategic environments and opening new avenues for research at the intersection of AI and game theory.
